\section*{Chapter \ref{ch:rs:portfolio-construction-i} --- Portfolio Construction I}

\begin{solution}{exo:rs:portfolio-construction-i:1}
$w_i = \alpha_i/(\gamma\sigma_i^2)$: $0.02/(5 \times 0.04) = 0.1$ and $0.01/(5 \times 0.01) = 0.2$. $\mathrm{IR} = \sqrt{0.02^2/0.04 + 0.01^2/0.01} = \sqrt{0.02} = 0.14$.
\end{solution}

\begin{solution}{exo:rs:portfolio-construction-i:2}
$84 \times 12 \times 0.0010 = 1.01$: about the book's whole capital every year, which is why its information ratio falls from 0.92 to 0.43.
\end{solution}

\begin{solution}{exo:rs:portfolio-construction-i:3}
$\sqrt{(1 + \mathrm{IR}^2/2)/T}$: $\sqrt{(1 + 1.77^2/2)/7.9} = 0.57$ and $\sqrt{(1 + 1.39^2/2)/7.9} = 0.50$. The two differ by less than one standard error.
\end{solution}

\begin{solution}{exo:rs:portfolio-construction-i:4}
With 126 observations of 100 returns, the sample covariance's smallest eigenvalues are far below the true ones (at the ratio 100/126 many are near zero). The
optimum $\Sigma^{-1}\alpha/\gamma$ divides by them: combinations that were nearly riskless in the sample get huge weights, long and short, and the gross explodes.
Their true risk is not small, so the realised volatility (207\%) dwarfs the forecast (36\%).
\end{solution}

\begin{solution}{exo:rs:portfolio-construction-i:5}
The forecast is a persistent true signal plus noise that is new each month. A book allowed to move only part of the way to each month's target holds an average of
recent targets, in which the noise partly cancels and the signal does not. Before costs, it trades a little less signal for much less noise.
\end{solution}

\begin{solution}{exo:rs:portfolio-construction-i:6}
$0.05 \times 10\,\mathrm{m}/1\,000\,\mathrm{m} = 0.05\%$ of capital.
\end{solution}

\begin{solution}{exo:rs:portfolio-construction-i:7}
The gross limit of 2 binds only if the book would otherwise exceed it. With the liquidity and turnover limits the gross averages 1.01, so the limit is slack every
month, its multiplier is zero, and so is its pull and its share (complementary slackness).
\end{solution}

\begin{solution}{exo:rs:portfolio-construction-i:8}
The ex-ante ratio measures the optimiser's own errors: with a six-month sample covariance of a hundred names the book's realised volatility was 207\% against 36\%
forecast, and its realised ratio 0.92 before costs and 0.43 after. Size on realised, cost-adjusted performance of a book built with a factor risk model and scaled
alphas, with its standard error.
\end{solution}

\begin{solution}{pb:rs:portfolio-construction-i:1}
\begin{enumerate}
\item $w^* = \Sigma^{-1}\alpha/\gamma$; $\mathrm{IR} = \sqrt{\alpha^\top\Sigma^{-1}\alpha}$.
\item The alpha is the simulation's stock-specific monthly expected return plus noise twice as large, standardised and scaled by $\mathrm{IC}\,\sigma_i$ with
$\mathrm{IC} = 0.05$; the risk model is chapter~24's at a monthly horizon. The truth plus noise isolates construction from forecasting.
\item Gross 61; forecast 36\%, realised 207\% volatility; ex-ante 4.25, realised 0.92 and 0.43 after costs; turnover 84 a month.
\item The smallest eigenvalues of the sample covariance belong to combinations of many names; the top two positions hold 9\% of the gross on average.
\item Forecast and realised volatility agree (13.7\% and 14.0\%), and the net information ratio is 1.39.
\item 1.64 to 1.59: the alpha is stock-specific, so the neutral book loses little.
\item The gross falls from 4.35 to 2.21 and 2.00, costs from 6.8\% to 3.0\% and 2.8\% a year.
\item 28 on average; the ex-ante ratio falls from 1.44 to 1.23.
\item Realised ratios over 7.9 years have standard errors near 0.5.
\item Turnover 0.53 a month, costs 0.6\% a year.
\item It averages monthly forecasts whose noise is independent from month to month.
\item The book no longer follows the noisy forecast; the correlation with it falls (0.97 to 0.58) while the correlation with the truth rises.
\item At the optimum $\alpha - \gamma\Sigma w^* = \sum_k g_k$, and $\alpha^\top\Sigma^{-1}\alpha - \gamma^2 w^{*\top}\Sigma w^* = \sum_k g_k^\top\Sigma^{-1}(2\alpha - G)$.
\item \textbf{Named result.} The turnover limit takes 44.0\% of the unconstrained ex-ante $\mathrm{IR}^2$, liquidity 39.4\%, the neutralities 7.0\%, the \omterm{def:rs:portfolio-construction-i:constraints}{name limits}
3.1\%, the gross limit and \omterm{def:rs:portfolio-construction-i:constraints}{dollar neutrality} nothing. After costs the turnover limit earned the most: the net ratio rises from 0.46 to 1.77.
\item Neither binds: the book's gross is about 1 once liquidity and turnover act, and \omterm{def:rs:portfolio-construction-i:constraints}{dollar neutrality} costs nothing with a stock-specific alpha.
\item The fully constrained book, reported with its forecast and realised volatility, its costs, the shadow prices and the decomposition, and the standard error of
its information ratio.
\item When the number of periods is many times the number of names, or as an input to shrinkage, never raw with more names than observations.
\item That they shrink the weights where estimation error is largest, so a constraint wrong on paper can reduce realised risk.
\item In factor form ($y = X^\top w$), with a sparse or splitting solver, warm starts and screening of names that cannot pay for their costs.
\item As much as its inputs are right: nothing, when its inputs are a sample covariance of more names than it has periods.
\end{enumerate}
\end{solution}

\begin{solution}{iq:rs:portfolio-construction-i:1}
They are error maximisers (Michaud): the optimiser favours assets with overestimated returns, underestimated variances and understated correlations, which are
estimation errors. Remedies: scale alphas by risk, use factor risk models and shrinkage, constrain, and penalise turnover.
\end{solution}

\begin{solution}{iq:rs:portfolio-construction-i:2}
The rate at which the objective would improve per unit of relaxation of a constraint. Use it to see which constraints bind and what they cost, to negotiate limits
with risk managers, and, as vectors, to decompose what each constraint takes from the forecast's information ratio.
\end{solution}

\begin{solution}{iq:rs:portfolio-construction-i:3}
The part of the alpha that lies along the neutralised exposures is lost; if the alpha is stock-specific (uncorrelated with the exposures) the cost is near zero
(here 1.64 to 1.59 ex ante). Jacobs, Levy and Starer showed such neutrality is in general suboptimal; it is chosen for risk and mandate reasons.
\end{solution}

\begin{solution}{iq:rs:portfolio-construction-i:4}
As a fraction of average daily traded value relative to capital, so the book can be exited in a few days; as capital grows the bounds tighten, more names bind
(28 of 100 here at \$1 billion), the book shifts to liquid names and the ex-ante ratio falls: capacity (chapter~28).
\end{solution}

\begin{solution}{iq:rs:portfolio-construction-i:5}
The correlation between the book's risk-adjusted weights and the risk-adjusted forecast (Clarke, de Silva and Thorley); it measures how much of the forecast reaches
the book. When the forecast is noisy, a book that deliberately does not follow it (a turnover limit) has a lower coefficient and a better realised ratio.
\end{solution}

\begin{solution}{iq:rs:portfolio-construction-i:6}
Factor form with exposures as variables; a sparse interior-point or operator-splitting solver; warm starts from yesterday's solution; precomputed factor
covariance, specific variances, exposures and liquidity bounds; screening of names whose alpha cannot pay for costs; checks on binding constraints and shadow
prices before trading.
\end{solution}
